\section{Část I: Společná matematika}

{\small Nižší priorita: ze společných okruhů (matematika + informatika) stačí 5/12.
Cíl je \textbf{breadth, ne hloubka} - u každého tématu umět definice a hlavní pojmy
(žlutá/zelená), aby z žádné přidělené otázky nebyla nula. Dlouhé důkazy nejsou potřeba.}

% ============================================================
\subsection{Základy diferenciálního a integrálního počtu}
\priorita{střední}{žlutá-zelená (1-2 body)}

\begin{tema}
Posloupnosti a jejich limity (aritmetika limit, věta o dvou policajtech); řady
(geometrická, harmonická); limita funkce a spojitost (nabývání mezihodnot a maxima);
derivace (l'Hospital, průběh funkce: extrémy, monotonie, konvexita; Taylorův polynom);
integrály (primitivní funkce, Riemannův/Newtonův integrál; aplikace: odhady součtů řad,
obsahy, objemy a povrchy rotačních útvarů, délka křivky).
\end{tema}

\ucivo
\begin{intuice}
\emph{Limita} = kam se hodnoty blíží. \emph{Derivace} = okamžitá rychlost změny (sklon
tečny). \emph{Integrál} = orientovaná plocha pod grafem. Derivace a integrál jsou navzájem inverzní
(základní věta analýzy).
\end{intuice}

\begin{odpoved}
\setlength{\parskip}{2pt}
\noindent\bod{1}\ \textbf{Minimum:} $\lim a_n=L \iff \forall\varepsilon>0\ \exists n_0\ \forall n\ge n_0:\ |a_n-L|<\varepsilon$.
$\lim_{x\to a}f(x)=L \iff \forall\varepsilon>0\ \exists\delta>0:\ 0<|x-a|<\delta \Rightarrow |f(x)-L|<\varepsilon$.
Spojitost v $a$: $\lim_{x\to a}f(x)=f(a)$; nespojitost v $a$ nastane, když $f(a)$ není definována, limita neexistuje nebo se nerovná $f(a)$.
Derivace $f'(a)=\lim_{h\to0}\frac{f(a+h)-f(a)}{h}$.\par
\noindent\bod{2}\ \textbf{+ k tomu:} \emph{Věta o dvou policajtech}: sevřená posloupnost má limitu
jako její meze. \emph{Geometrická řada} $\sum q^n=\frac{1}{1-q}$ pro $|q|<1$, harmonická
$\sum\frac1n$ diverguje. \emph{l'Hospital}: pro typ $\frac00$ nebo $\frac\infty\infty$
lze za podmínek věty nahradit podílem derivací. Průběh: $f'>0$ roste, $f''>0$ konvexní,
kandidáty na extrém dávají body s $f'=0$ nebo kde derivace není definovaná. \emph{Newton-Leibniz}:
\emph{primitivní funkce} na intervalu $I$ splňuje $F'=f$ na $I$ a $\int_a^b f = F(b)-F(a)$.
\emph{Riemannův integrál} je limita součtů $\sum_i f(\xi_i)(x_i-x_{i-1})$, pokud při zjemnění dělení existuje a nezávisí na volbě $\xi_i$.\par
\noindent\bod{3}\ \textbf{Bonus:} \emph{Taylor} $f(x)\approx\sum\frac{f^{(k)}(a)}{k!}(x-a)^k$.
Aplikace integrálu: objem rotačního tělesa $\pi\int f(x)^2\,dx$, délka křivky
$\int\sqrt{1+f'(x)^2}\,dx$, povrch rotační plochy
$2\pi\int f(x)\sqrt{1+f'(x)^2}\,dx$. Integrální kritérium: pro nezápornou
nerostoucí $f$ konverguje $\sum f(k)$ právě když je nevlastní integrál konečný.
\end{odpoved}
\begin{recall}
\begin{itemize}[leftmargin=1.2em,itemsep=0pt,topsep=1pt]
  \item Kdy konverguje geometrická řada a jaký má součet?
  \item Co poznám z $f'$ a co z $f''$?
  \item Jak souvisí Riemannův integrál s primitivní funkcí?
\end{itemize}
\end{recall}
\past{Harmonická řada diverguje, i když členy jdou k nule. l'Hospital jen na neurčité výrazy.}

% ============================================================
\subsection{Algebra a lineární algebra}
\priorita{střední}{žlutá-zelená (1-2 body)}

\begin{tema}
Algebraické struktury (grupy, permutace, tělesa včetně konečných); soustavy lineárních
rovnic (Gaussova a Gauss-Jordanova eliminace); matice (hodnost, regulární a inverzní);
vektorové prostory (lineární nezávislost, báze, dimenze, Steinitzova věta, podprostory);
lineární zobrazení (obraz, jádro, isomorfismus); skalární součin (norma, Cauchy-Schwarz,
ortonormální báze, Gram-Schmidt, ortogonální projekce a matice); determinanty (Laplaceův
rozvoj, geometrie); vlastní čísla a vektory (charakteristický polynom, diagonalizace,
spektrální rozklad, symetrické matice); pozitivně (semi)definitní matice (Choleského rozklad).
\end{tema}

\ucivo
\begin{intuice}
Lineární algebra studuje \emph{vektorové prostory} a \emph{lineární zobrazení} (= násobení
maticí). Hodně pojmů jsou jen různé pohledy na totéž: hodnost, dimenze obrazu, počet pivotů.
\end{intuice}

\begin{odpoved}
\setlength{\parskip}{2pt}
\noindent\bod{1}\ \textbf{Minimum:} \emph{Hodnost} matice je dimenze prostoru generovaného jejími sloupci, ekvivalentně počet pivotů.
\emph{Báze} vektorového prostoru je lineárně nezávislá generující množina; \emph{dimenze} je počet prvků libovolné báze.
\emph{Vlastní číslo a vektor}: $Ax=\lambda x$ pro $x\neq0$. Matice je regulární $\iff$ má inverzi $\iff$ $\det\neq0$ $\iff$ plná hodnost.\par
\noindent\bod{2}\ \textbf{+ k tomu:} \emph{Gaussova eliminace} řeší soustavy i počítá hodnost/inverzi.
\emph{Jádro} = řešení $Ax=0$; platí $\dim\ker + \mathrm{rank} = n$. Vlastní čísla jsou kořeny char.\ polynomu $\det(A-\lambda I)=0$.
\emph{Podobné matice}: $B=S^{-1}AS$ pro regulární $S$; zachovávají hodnost, determinant, vlastní čísla, stopu a charakteristický polynom.
\emph{Pozitivně definitní} reálná symetrická matice splňuje $x^\top Ax>0$ pro všechna $x\neq0$.
\emph{Skalární součin} dává normu a kolmost;
\emph{Gram-Schmidt} vyrobí ortonormální bázi. Souřadnice vektoru vůči ortonormální bázi jsou
\emph{Fourierovy koeficienty} $\langle u|v_i\rangle$ (a \emph{ortogonální projekce} na podprostor
je součet $\sum_i \langle u|v_i\rangle v_i$). Pythagorova věta a Cauchy-Schwarzova nerovnost $|\langle u|v\rangle|\le\|u\|\|v\|$.\par
\noindent\bod{3}\ \textbf{Bonus:} \emph{Diagonalizace} $A=S\Lambda S^{-1}$ (sloupce $S$ = vlastní
vektory). Symetrické matice mají reálná vl.\ čísla a ortogonální diagonalizaci $A=Q\Lambda Q^\top$.
Reálná symetrická matice je \emph{pozitivně definitní} ($x^\top Ax>0$ pro $x\neq0$)
$\iff$ má kladná vl.\ čísla $\iff$ existuje Choleského rozklad $A=LL^\top$.
Grupa má asociativní operaci, neutrální a inverzní prvky; těleso má sčítání a násobení
kompatibilní distributivitou, konečné těleso má řád mocniny prvočísla. Permutace se skládá
z cyklů, rozkládá na transpozice a má znaménko. Souřadnice a matice lineárního zobrazení
závisí na bázi; matice přechodu je regulární. Ortogonální matice zachovává skalární součin
a délky. Determinant měří orientovaný objem, je multiplikativní a stopa je součet vlastních čísel.
\end{odpoved}
\begin{recall}
\begin{itemize}[leftmargin=1.2em,itemsep=0pt,topsep=1pt]
  \item Co znamená, že je matice regulární (vyjmenuj ekvivalentní podmínky)?
  \item Jak najdu vlastní čísla a co je vlastní vektor?
  \item Co je pozitivně definitní matice a jak souvisí s vlastními čísly?
\end{itemize}
\end{recall}
\past{Vlastní čísla $=$ kořeny char.\ polynomu, ne prvky na diagonále (pokud matice není trojúhelníková).}

% ============================================================
\subsection{Diskrétní matematika}
\priorita{střední}{žlutá-zelená (1-2 body)}

\begin{tema}
Relace a jejich vlastnosti (reflexivita, symetrie, antisymetrie, tranzitivita); ekvivalence
a rozkladové třídy; částečná uspořádání (řetězec, antiřetězec, výška a šířka a věta o jejich
vztahu); funkce (prostá, na, bijekce; počty funkcí mezi konečnými množinami); permutace;
kombinační čísla a binomická věta; princip inkluze a exkluze (důkaz a aplikace); Hallova
věta o systému různých reprezentantů (vztah k párování v bipartitním grafu, polynomiální algoritmus).
\end{tema}

\ucivo
\begin{intuice}
Hodně z toho je „počítání možností“ a struktura na konečných množinách. \emph{Inkluze-exkluze}
opravuje dvojí započítání; \emph{Hallova věta} říká, kdy lze každého spárovat.
\end{intuice}

\begin{odpoved}
\setlength{\parskip}{2pt}
\noindent\bod{1}\ \textbf{Minimum:} \emph{Ekvivalence} je relace reflexivní, symetrická a tranzitivní; její třídy tvoří rozklad množiny.
\emph{Částečné uspořádání} je relace reflexivní, antisymetrická a tranzitivní.
\emph{Bijekce} = prostá i na. Počet funkcí $A\to B$ je $|B|^{|A|}$.\par
\noindent\bod{2}\ \textbf{+ k tomu:} V částečném uspořádání je \emph{řetězec} po dvou porovnatelná podmnožina a \emph{antiřetězec} po dvou neporovnatelná podmnožina;
\emph{Dilworthova věta} dává vztah šířky a pokrytí řetězci.
\emph{Inkluze-exkluze}: $|A\cup B|=|A|+|B|-|A\cap B|$ (obecně se střídají
znaménka). \emph{Binomická věta}: $(a+b)^n=\sum\binom nk a^k b^{n-k}$ (kombinační čísla =
Pascalův trojúhelník).\par
\noindent\bod{3}\ \textbf{Bonus:} \emph{Hallova věta}: bipartitní graf má párování pokrývající $A$
$\iff$ pro každou $X\subseteq A$ platí $|N(X)|\ge|X|$ (Hallova podmínka); související
\emph{Kőnigova věta}: max.\ párování = min.\ vrcholové pokrytí v bipartitním grafu. Aplikace inkluze-exkluze:
počet surjekcí, Eulerova funkce, problém šatnářky. Počet prostých funkcí $N\to M$ je
$m^{\underline n}$ a bijekcí na $n$ prvcích je $n!$. V konečné neprázdné ČUM existují
minimální a maximální prvky, výška krát šířka je aspoň $|X|$.
\end{odpoved}
\begin{recall}
\begin{itemize}[leftmargin=1.2em,itemsep=0pt,topsep=1pt]
  \item Jaké tři vlastnosti dělají z relace ekvivalenci?
  \item Jak zní Hallova podmínka?
  \item Kolik je funkcí (a kolik bijekcí) mezi konečnými množinami?
\end{itemize}
\end{recall}
\past{Antisymetrie neznamená „není symetrická“. Bijekce mezi $A,B$ existuje jen při $|A|=|B|$.}

% ============================================================
\subsection{Teorie grafů}
\priorita{střední}{zelená (2 body), vděčné téma}

\begin{tema}
Základní pojmy (izomorfismus, stupeň vrcholu, doplněk, bipartitní graf); příklady (úplný,
cesty, kružnice); souvislost, komponenty, vzdálenost; stromy (charakteristiky); rovinné grafy
(Eulerova formule a max.\ počet hran); barevnost (vztah barevnosti a klikovosti); hranová a
vrcholová souvislost (Mengerova věta); orientované grafy (silná a slabá souvislost); toky v
sítích (definice, existence max.\ toku, Ford-Fulkerson pro celočíselné kapacity).
\end{tema}

\ucivo
\begin{intuice}
Grafy = vrcholy a hrany. Mnoho úloh = „dá se projít / spojit / obarvit / protlačit tok?“.
\end{intuice}

\begin{center}
\begin{tikzpicture}[>=Stealth, node distance=11mm and 17mm,
   every node/.style={circle,draw,minimum size=6mm,font=\small,inner sep=0pt}]
  \node (s) {$s$};
  \node (a) [above right=of s] {$a$};
  \node (b) [below right=of s] {$b$};
  \node (t) [below right=of a] {$t$};
  \tikzset{every edge/.style={draw,->,font=\scriptsize}}
  \draw (s) edge node[above left]{$3$} (a);
  \draw (s) edge node[below left]{$2$} (b);
  \draw (a) edge node[left]{$1$} (b);
  \draw (a) edge node[above right]{$2$} (t);
  \draw (b) edge node[below right]{$3$} (t);
\end{tikzpicture}\\[-2pt]
{\scriptsize Síť s kapacitami. Max tok = min řez (zde 5).}
\end{center}

\begin{odpoved}
\setlength{\parskip}{2pt}
\noindent\bod{1}\ \textbf{Minimum:} \emph{Strom} je souvislý graf bez kružnic, ekvivalentně souvislý graf na $n$ vrcholech s $n-1$ hranami.
\emph{Bipartitní graf} má rozklad $V=A\dot\cup B$ tak, že každá hrana vede mezi $A$ a $B$, ekvivalentně nemá lichou kružnici.
\emph{Barevnost} $\chi(G)$ je nejmenší počet barev vrcholů tak, aby konce každé hrany měly různé barvy.\par
\noindent\bod{2}\ \textbf{+ k tomu:} \emph{Silně souvislý} orientovaný graf má orientovanou cestu mezi každou uspořádanou dvojicí vrcholů; \emph{kondenzace} je DAG silných komponent.
\emph{Eulerovský tah} projde každou hranu právě jednou, eulerovský graf má uzavřený eulerovský tah. \emph{Rovinný graf} lze nakreslit bez křížení hran;
pro souvislý rovinný graf platí Eulerova formule $v-e+f=2$ a pro jednoduchý rovinný graf s $v\ge3$ plyne max.\ $3v-6$ hran.
\emph{Tok} v síti respektuje kapacity a zachování toku; hodnota max.\ toku se rovná kapacitě min.\ řezu.\par
\noindent\bod{3}\ \textbf{Bonus:} \emph{Mengerova věta}: pro dva různé vrcholy (u vrcholové verze typicky nesousední) je maximum hranově
(resp.\ vnitřně vrcholově) disjunktních cest rovno velikosti minimálního hranového (resp.\ vrcholového) separátoru.
Slabě souvislý orientovaný graf má souvislý podkladový graf.
Eulerovský graf je souvislý a všechny stupně jsou sudé; v orientovaném grafu potřebuje
vyvážené vstupní a výstupní stupně a souvislý podklad po odstranění izolovaných vrcholů.
\emph{Ford-Fulkerson} opakuje zlepšující cesty v reziduální síti. Kuratowski: nerovinný graf obsahuje dělení $K_5$ nebo
$K_{3,3}$; rovinný graf bez trojúhelníků má nejvýše $2v-4$ hran.
\end{odpoved}
\begin{recall}
\begin{itemize}[leftmargin=1.2em,itemsep=0pt,topsep=1pt]
  \item Kolik hran má strom a jak poznám bipartitní graf?
  \item Co říká Eulerova formule a kolik nejvíc hran má rovinný graf?
  \item Jak funguje Ford-Fulkerson a co je max-flow = min-cut?
\end{itemize}
\end{recall}
\past{Eulerova formule platí pro souvislý rovinný graf; odhad $3v-6$ navíc předpokládá jednoduchost a $v\ge3$. Bipartitní $\iff$ bez liché kružnice.}

% ============================================================
\subsection{Pravděpodobnost a statistika}
\priorita{střední}{žlutá-zelená (1-2 body)}

\begin{tema}
Pravděpodobnostní prostor, jevy, nezávislost, podmíněná pravděpodobnost, Bayesův vzorec;
náhodné veličiny (diskrétní i spojité), distribuční funkce, hustota/pravděpodobnostní funkce;
střední hodnota (linearita, Markovova nerovnost); rozptyl (vzorec pro součet); konkrétní
rozdělení (geometrické, binomické, Poissonovo, normální, exponenciální); limitní věty (zákon
velkých čísel, centrální limitní věta); bodové a intervalové odhady; testování hypotéz (chyby
1.\ a 2.\ druhu, hladina významnosti).
\end{tema}

\ucivo
\begin{intuice}
Pravděpodobnost popisuje nejistotu. \emph{Střední hodnota} = dlouhodobý průměr, \emph{rozptyl} =
míra rozkolísanosti. \emph{CLV} říká, proč se všude objevuje normální rozdělení.
\end{intuice}

\begin{center}
\begin{tikzpicture}
\begin{axis}[width=6.6cm,height=3cm,axis lines=left,domain=-3.4:3.4,samples=80,
  xtick=\empty,ytick=\empty,clip=false]
  \addplot[thick,defblue]{exp(-x^2/2)};
\end{axis}
\end{tikzpicture}\\[-4pt]
{\scriptsize Normální rozdělení (Gaussova křivka): limita součtu mnoha nezávislých vlivů (CLV).}
\end{center}

\begin{odpoved}
\setlength{\parskip}{2pt}
\noindent\bod{1}\ \textbf{Minimum:} \emph{Pravděpodobnostní prostor} je trojice $(\Omega,\mathcal F,P)$, kde $\mathcal F$ je $\sigma$-algebra jevů a $P(\Omega)=1$ je spočetně aditivní.
\emph{Podmíněná pravděpodobnost} pro $P(B)>0$ je $P(A|B)=\frac{P(A\cap B)}{P(B)}$.
\emph{Bayes} $P(A|B)=\frac{P(B|A)P(A)}{P(B)}$. \emph{Střední hodnota} $E[X]=\sum x_i p_i$
(resp.\ $\int x f(x)$).\par
\noindent\bod{2}\ \textbf{+ k tomu:} \emph{Linearita} $E[X+Y]=E[X]+E[Y]$ (vždy). \emph{Rozptyl}
$\mathrm{var}(X)=E[X^2]-E[X]^2$; obecně
$\mathrm{var}(X+Y)=\mathrm{var}X+\mathrm{var}Y+2\mathrm{cov}(X,Y)$, pro nezávislé kovariance zmizí.
\emph{Markovova nerovnost}: pro $X\ge0$ a $a>0$ platí $P(X\ge a)\le E[X]/a$.\par
\noindent\bod{3}\ \textbf{Bonus:} \emph{Zákon velkých čísel}: průměr $\to E[X]$. \emph{CLV}: pro nezávislé stejně rozdělené veličiny se střední hodnotou $\mu$ a rozptylem $\sigma^2>0$ platí
$\frac{\sum_{i=1}^n X_i-n\mu}{\sigma\sqrt n}\Rightarrow N(0,1)$.
\emph{Test hypotéz}: chyba 1.\ druhu =
zamítnu platnou $H_0$ (hladina $\alpha$), 2.\ druhu = nezamítnu neplatnou.
Konkrétní rozdělení: Bernoulli, geometrické, binomické, Poissonovo, uniformní,
exponenciální a normální. Bodové odhady řeší nestrannost, konzistenci, bias a MSE;
tvoří se metodou momentů nebo MLE. Konfidenční interval má spolehlivost $1-\alpha$,
p-hodnota je nejmenší hladina, na níž zamítáme.
\end{odpoved}
\begin{recall}
\begin{itemize}[leftmargin=1.2em,itemsep=0pt,topsep=1pt]
  \item Jak zní Bayesův vzorec?
  \item Kdy platí $\mathrm{var}(X+Y)=\mathrm{var}X+\mathrm{var}Y$?
  \item Co říká centrální limitní věta?
\end{itemize}
\end{recall}
\past{Linearita střední hodnoty platí vždy, ale součet rozptylů jen pro nezávislé veličiny.}

% ============================================================
\subsection{Logika}
\priorita{střední}{žlutá-zelená (1-2 body)}

\begin{tema}
Syntaxe (výroková a predikátová logika, normální tvary CNF/DNF/PNF, převody, použití pro SAT
a rezoluci); sémantika (model teorie, splnitelnost, tautologie, důsledek); extenze teorií
(konzervativnost, skolemizace); dokazatelnost (formální důkaz, dokazovací systémy: tablo,
rezoluce, příp.\ Hilbertovský kalkul); věty o kompaktnosti a úplnosti; rozhodnutelnost.
\end{tema}

\ucivo
\begin{intuice}
\emph{Syntaxe} = forma formulí, \emph{sémantika} = jejich pravdivost v modelech. \emph{Dokazatelnost}
(syntaktická) a \emph{platnost} (sémantická) spojuje věta o úplnosti.
\end{intuice}

\begin{odpoved}
\setlength{\parskip}{2pt}
\noindent\bod{1}\ \textbf{Minimum:} \emph{Tautologie} = pravdivá ve všech ohodnoceních;
\emph{splnitelná} = pravdivá aspoň v jednom. $T\models\varphi$: $\varphi$ platí v každém modelu $T$.
\emph{CNF} = konjunkce klauzulí, \emph{DNF} = disjunkce konjunkcí.\par
\noindent\bod{2}\ \textbf{+ k tomu:} \emph{Rezoluce} (na CNF): z klauzulí $C\lor\ell$ a $D\lor\neg\ell$
odvodí $C\lor D$; spor = prázdná klauzule (důkaz nesplnitelnosti). \emph{Tablo} = systematické
hledání modelu. \emph{Věta o úplnosti}: $T\models\varphi \iff T\vdash\varphi$.\par
\noindent\bod{3}\ \textbf{Bonus:} \emph{Kompaktnost}: teorie má model $\iff$ každá konečná část má
model. \emph{Skolemizace} odstraní existenční kvantifikátory. \emph{Rozhodnutelnost}: výroková
logika je rozhodnutelná (SAT), predikátová obecně ne; \emph{rekurzivně axiomatizovaná kompletní}
teorie je rozhodnutelná, ale např.\ Peanova aritmetika kompletní není.
PNF má kvantifikátory v prefixu; Horn-SAT je lineární přes jednotkovou propagaci.
Extenze je konzervativní, když nepřidá nové důsledky ve starém jazyce. Tablo je
sporné, když se uzavřou všechny větve; korektnost a úplnost spojují důkaz s platností.
Gödel: dost silná bezesporná rekurzivně axiomatizovaná aritmetika je nekompletní.
\end{odpoved}
\begin{recall}
\begin{itemize}[leftmargin=1.2em,itemsep=0pt,topsep=1pt]
  \item Rozdíl mezi tautologií, splnitelností a důsledkem?
  \item Jak funguje rezoluční pravidlo a co dokazuje prázdná klauzule?
  \item Co říká věta o úplnosti?
\end{itemize}
\end{recall}
\past{Nezaměňuj „splnitelná“ (aspoň jeden model) a „tautologie“ (všechny modely). Rezoluce dokazuje nesplnitelnost (hledá spor).}
